\documentclass{article}
\usepackage{graphicx}
\usepackage[numbers]{natbib}
\usepackage{float}
\usepackage{hyperref}
\usepackage{comment}
\usepackage{ragged2e}

\begin{titlepage}
    \title{Járművek trajektóriájinak előrejelzése machine learning módszerrel}
    \author{Péter Bence\\Mérnökinformatika BSc}
    \date{2023 Február}
\end{titlepage}

\begin{document}
    \bibliographystyle{plainnat}
    \begin{titlepage}
        \maketitle
        \begin{figure}[H]\centering
            \includegraphics[width=1\columnwidth]{sze_givk_logo.png}
        \end{figure}
    \end{titlepage}
    \begin{abstract}
        Az ITS (intelligent transportation system) egyre nagyobb teret 
        hódít napjainkban és rengeteg különböző területen alkalmazzák 
        ezeket a rendszereket. A közlekedési csomópontok elemzése egy 
        frekventált terület az ITS alkalmazásában.
        Célunk, gépi látás és gépi tanulás felhasználásá-val, közlekedési
        csomópontok elemzésének automatizálása és felgyorsítása. A kutatásban
        lefektetett alapgondolatokat, kifejlesztett keretrendszert és 
        a felmerülő probélmák megoldásait, a gyakorlatban balesetek megelőzésére,
        renitens viselkedések kiszűrésére és forgalomirányító renszerek támogatásá-ra
        lehet használni.  
        A kutatásban egy trajektória osztályozó módszert ismertetünk, amely objektumdetektálás 
        és objektumkövetés segítségével elemezi a közle-kedési 
        csomópontokban elhaladó járművek mozgását. A mozgásuk alapján 
        automatikusan klaszterezi a trajektóriákat, majd gépi tanulás 
        segítségével predikciót ad az újonnan belépő járművek kilépési 
        pontjára. A módszerhez 5 különböző közlekedési csomópontban 
        készített saját videó adatbázisunkat használtuk fel.
        A tesztelt klaszterezési mód-szerek közül (OPTICS, BIRCH, KMeans, DBSCAN)
        az OPTICS algoritmus bizonyult legjobbnak trajektórák klaszterezésére.
        Összehasonlí-tottunk több különböző klasszifikációs módszert 
        a legpontosabb predikció eléréséhez, amelyek: KNN, SVM, GP, DT, 
        GNB, MLP, SGD. A tanul-mányban bemutatott eljárások közül az 
        SVM adta a legpontosabb 90\%-os eredményt.
        Ezt a pontosságot valós idejű futás közben 5 fps mellett érte el.
        Ebből azt a következtetést lehet levonni, hogy jobb sebesség elérése
        érdekében, vagy a feature vectorok dimenzióját kell csökkenteni, vagy
        érdemes neurális hálót alkalmazni a klasszifikációhoz.
        A forráskód megtalálható ebben a github repositoriban 
        \url{https://www.github.com/Pecneb/computer_vision_research}
        
        %Ebben a kutatásban gépi tanulási módszerrel törekszünk megoldani
        %járművek trajektóriájának előrejelzését komplex forgalmi helyzetekben.
        %Madártávlatból felvett felvételekből nyertünk ki járművek trajektóriájit 
        %YOLO konvolúciós neurális háló és DeepSORT objektum követő algoritmus segítségével. 
        %Különböző klaszterezése algoritmusokat 
        %tesztel-tünk, többfajta feature vectorral, hogy minél pontosabban kategorizálni 
        %tudjuk a járművek mozgását. Ezek kategóriák felhasználásával klasszifikációs modelleket 
        %tanítottunk be, több fajta feature vectorral, amik-nek pontosságát megmértük 
        %és összehasonlítottuk. Kutatásunk során a teljes tanítási és tesztelési folyamat 
        %gyorsítására és automatizálására nagy hang-súlyt fektettünk, ennek megoldására egy 
        %keretrendszert fejlesztettünk ki. A kutatásból azt a tanulságot tudjuk leszűrni,
        %hogy a klaszterezési folyamatot nehéz automatizálni, mivel a klaszterezéshez
        %használt modellek hiperparamétereit nem lehet univerzálisan minden ke-reszteződésre 
        %felhasználni. A klasszifikációhoz használt modellek nem csak pontosnak,
        %hanem gyorsnak is kell lennie, ha valós időben akarjuk használni.
        %A detektálásnál felhasznált DeepSORT objektum követő algoritmus
        %hiperparamétereit is gondosan meg kell választani, hogy ne legyenek anomáliák a tanító
        %adathalmazban. A forráskód megtalálható ebben a github repositoriban \url{https://www.github.com/Pecneb/computer_vision_research}
    \end{abstract}
    %\tableofcontents

    %\bibliography{biblio}
\end{document}